% Paste this environment into the IEEEtran source.
% The preamble must already \input the sibling iscafig.tex.
\begin{figure*}[t]
\centering
\begin{tikzpicture}[iscafig]
  % kind bits
  \node[lbl, anchor=west] at (0,0.42)
    {\texttt{MUT\_COLLEC} payload: address in [31:2], kind in [1:0]};
  \node[fldimm, minimum width=3.85cm, minimum height=6.4mm, anchor=west] at (0,0) {00\quad list};
  \node[fldimm, minimum width=3.85cm, minimum height=6.4mm, anchor=west] at (4.05,0) {01\quad dict};
  \node[fldimm, minimum width=3.85cm, minimum height=6.4mm, anchor=west] at (8.10,0) {10\quad set};
  \node[fldimm, minimum width=4.15cm, minimum height=6.4mm, anchor=west] at (12.15,0) {11\quad bytearray};

  % list
  \begin{scope}[yshift=-1.15cm]
    \node[lbl, anchor=west] at (0,0.38) {\textbf{List}, three header words};
    \rec{0}{+0}{fldimm}{6.4cm}{length}
    \rec{-0.55}{+4}{fldimm}{6.4cm}{capacity, minimum 4}
    \rec{-1.10}{+8}{fldmem}{6.4cm}{pointer to item 0}
    \node[fldcore, anchor=west, minimum width=6.4cm] (items) at (0,-1.85)
      {item $i$, one tagged word};
    \draw[thinarr] ($(0,-1.10)+(6.55,0)$) -- ++(0.35,0) |- (items.east);
    \node[lbl, anchor=west, align=left] at (0,-2.45)
      {A full list doubles the item array and\\
       updates the pointer. The header does not move.};
  \end{scope}

  % bytearray
  \begin{scope}[xshift=8.7cm, yshift=-1.15cm]
    \node[lbl, anchor=west] at (0,0.38) {\textbf{Bytearray}, three header words};
    \rec{0}{+0}{fldimm}{6.4cm}{length, in bytes}
    \rec{-0.55}{+4}{fldimm}{6.4cm}{capacity, in bytes}
    \rec{-1.10}{+8}{fldmem}{6.4cm}{pointer to the byte buffer}
    \node[fldimm, anchor=west, minimum width=6.4cm] at (0,-1.85)
      {raw bytes, not scanned};
    \node[lbl, anchor=west, align=left] at (0,-2.45)
      {Same indexing shape as a list.\\
       The buffer holds bytes, not tagged values.};
  \end{scope}

  % dict
  \begin{scope}[yshift=-4.55cm]
    \node[lbl, anchor=west] at (0,0.38) {\textbf{Dict}, four header words};
    \rec{0}{+0}{fldhot}{6.4cm}{used\,[30:0], object bit\,[31]}
    \rec{-0.55}{+4}{fldimm}{6.4cm}{allocated slots}
    \rec{-1.10}{+8}{fldmem}{6.4cm}{pointer to entry 0}
    \rec{-1.65}{+12}{fldmem}{6.4cm}{types table, or 0}
    \node[fldcore, anchor=west, minimum width=3.1cm] (k) at (0,-2.40) {key};
    \node[fldcore, anchor=west, minimum width=3.1cm] (v) at (3.3,-2.40) {value};
    \node[lbl, anchor=west] at (0,-2.95)
      {empty $=$ \texttt{CONTROL}/0 \quad tombstone $=$ \texttt{CONTROL}/3};
  \end{scope}

  % set
  \begin{scope}[xshift=8.7cm, yshift=-4.55cm]
    \node[lbl, anchor=west] at (0,0.38) {\textbf{Set}, same header as a dict};
    \rec{0}{+0}{fldhot}{6.4cm}{used\,[30:0], object bit\,[31]}
    \rec{-0.55}{+4}{fldimm}{6.4cm}{allocated slots}
    \rec{-1.10}{+8}{fldmem}{6.4cm}{pointer to slot 0}
    \rec{-1.65}{+12}{fldmem}{6.4cm}{types table, or 0}
    \node[fldcore, anchor=west, minimum width=6.4cm] at (0,-2.40)
      {one key per slot};
    \node[lbl, anchor=west] at (0,-2.95)
      {tombstones are not copied on a resize};
  \end{scope}
\end{tikzpicture}
\caption{Mutable collections.
The low two bits of a \texttt{MUT\_COLLEC} payload select the kind, and the address clears them.
A list stores length and capacity in the object and keeps items in a separate array.
Dict and set share a four-word header.
Bit 31 of the first word stays set once an \texttt{OBJECT} has been inserted, and the types-table pointer is then live.
Green boxes are pointers.
An empty slot is \texttt{CONTROL} kind 0.
A deleted slot is \texttt{CONTROL} kind 3.}
\label{fig:collections}
\end{figure*}
